\section{Muối và phản ứng trong dung dịch}

\subsection{Muối và tính tan}

Muối là hợp chất gồm cation kim loại hoặc \ce{NH4+} và anion gốc acid. Khi giải bài phản ứng trong dung dịch, cần biết chất nào tan để xác định các ion thực sự có mặt trong dung dịch.

Không cần học thuộc toàn bộ bảng tính tan. Hãy tra bảng khi cần và ghi nhớ dần các chất xuất hiện thường xuyên.

\begin{ghinho}{Dấu hiệu giúp nhận ra phản ứng}
Hai chất trong dung dịch thường phản ứng khi tạo được ít nhất một trong các sản phẩm sau:
\begin{itemize}
\item chất kết tủa;
\item chất khí;
\item nước hoặc chất điện li yếu.
\end{itemize}
\end{ghinho}

\subsection{Các kiểu phản ứng thường gặp của muối}

\textbf{Muối + kim loại:}
\[
\ce{CuSO4 + Fe -> FeSO4 + Cu}.
\]

\textbf{Muối + acid:}
\[
\ce{Na2CO3 + 2HCl -> 2NaCl + CO2 ^ + H2O}.
\]

\textbf{Muối + base:}
\[
\ce{CuSO4 + 2NaOH -> Cu(OH)2 v + Na2SO4}.
\]

\textbf{Muối + muối:}
\[
\ce{BaCl2 + Na2SO4 -> BaSO4 v + 2NaCl}.
\]

\subsection{Cách dự đoán hiện tượng}

\begin{ghinho}{Quy trình}
\begin{enumerate}
\item Viết công thức các chất ban đầu.
\item Dự đoán sản phẩm trao đổi hoặc sản phẩm của phản ứng thế.
\item Dùng bảng tính tan và tính chất chất khí để kiểm tra phản ứng có xảy ra không.
\item Từ sản phẩm suy ra hiện tượng: kết tủa, khí, chất rắn bám, đổi màu dung dịch, chất rắn tan dần,...
\end{enumerate}
\end{ghinho}

\begin{vidu}
Trộn dung dịch \ce{BaCl2} với dung dịch \ce{Na2SO4}. Nêu hiện tượng và viết phương trình hóa học.
\end{vidu}
\begin{loigiai}
Hai muối trao đổi ion, tạo \ce{BaSO4} không tan:
\[
\ce{BaCl2 + Na2SO4 -> BaSO4 v + 2NaCl}.
\]
Xuất hiện kết tủa trắng.

\textbf{Vậy,} phản ứng tạo kết tủa trắng \ce{BaSO4}.
\end{loigiai}

\subsection{Carbonate và bicarbonate}

Carbonate và bicarbonate xuất hiện nhiều trong các bài về đá vôi, nước cứng và xử lí cặn.

\[
\ce{CaCO3 + 2HCl -> CaCl2 + CO2 ^ + H2O},
\]
\[
\ce{Ca(HCO3)2 ->[t^\circ] CaCO3 v + CO2 ^ + H2O}.
\]

Khi đun nước cứng tạm thời, bicarbonate bị phân hủy tạo cặn carbonate. Acid có thể hòa tan cặn carbonate và giải phóng \ce{CO2}.

\begin{tuluyen}
Nêu hiện tượng và viết phương trình khi cho dung dịch \ce{Na2CO3} vào dung dịch \ce{HCl} dư.
\dapso{Có khí không màu thoát ra; \ce{Na2CO3 + 2HCl -> 2NaCl + CO2 ^ + H2O}.}
\begin{loigiaiBT}
Carbonate phản ứng với acid tạo muối, carbon dioxide và nước:
\[
\ce{Na2CO3 + 2HCl -> 2NaCl + CO2 ^ + H2O}.
\]
\ce{CO2} thoát ra làm dung dịch sủi bọt khí.

\textbf{Vậy,} hiện tượng là có khí không màu thoát ra.
\end{loigiaiBT}
\end{tuluyen}
